\chapter{Cheat Sheet Summary}
This chapter is the print-and-review version of the book. Later chapters explain
where the formulas come from and connect them to state-space and
frequency-domain thinking.

\textbf{Reading these symbols.} Lowercase letters (\(v,i,q\)) are instantaneous
time-domain values; uppercase (\(V,I\)) are RMS phasors or steady DC values. The
impedance, time-constant, resonance, decibel, filter, and transmission-line
quantities below are the exam-essential ones. A handful of symbols---\(\zeta\)
(damping ratio), \(A\) and \(\beta\) (amplifier forward gain and feedback
fraction), \(L\) (loop gain), and \(\mathcal{H}\) (total stored energy)---appear
only to connect the formulas to control theory and are not quantities the exam
asks you to compute.

\section{Quantities, Units, and Symbols}
\begin{tabularx}{\textwidth}{@{}L{0.19\textwidth}L{0.13\textwidth}L{0.16\textwidth}Y@{}}
\toprule
Quantity & Symbol & SI unit & Meaning \\
\midrule
Charge & \(q,Q\) & \si{\coulomb} & Stored electric quantity. \\
Current & \(i,I\) & \si{\ampere} & Rate of charge flow, \(i=\dd q/\dd t\). \\
Voltage & \(v,V\) & \si{\volt} & Energy per unit charge. \\
Power & \(p,P\) & \si{\watt} & Rate of energy flow. \\
Energy & \(E,\mathcal{H}\) & \si{\joule} & Stored or dissipated energy; \(\mathcal{H}\) is total stored energy. \\
Resistance & \(R\) & \si{\ohm} & Opposition to current. \\
Conductance & \(G\) & \si{\siemens} & Reciprocal resistance, \(G=1/R\). \\
Capacitance & \(C\) & \si{\farad} & Charge stored per volt. \\
Inductance & \(L\) & \si{\henry} & Voltage-time per ampere. \\
Frequency & \(f\) & \si{\hertz} & Cycles per second. \\
Angular frequency & \(\omega\) & \si{\radian\per\second} & Radians per second, \(\omega=2\pi f\). \\
\bottomrule
\end{tabularx}

\section{Core Circuit Formulas}
\begin{longtable}{@{}L{0.25\textwidth}L{0.32\textwidth}L{0.34\textwidth}@{}}
\toprule
Topic & Formula & Assumption or memory cue \\
\midrule
\endhead
Ohm's law & \(V=IR\) & Linear resistor under passive sign convention. \\
Power & \(P=VI=I^2R=V^2/R\) & Use RMS values for sinusoidal average power in a resistor. \\
Charge/current & \(i=\dd q/\dd t\) & Current is rate of charge flow. \\
Capacitor law & \(i=C\,\dd v/\dd t\) & Capacitor voltage cannot jump under finite current. \\
Inductor law & \(v=L\,\dd i/\dd t\) & Inductor current cannot jump under finite voltage. \\
Capacitor energy & \(E_C=\frac{1}{2}CV^2\) & Electric-field storage. \\
Inductor energy & \(E_L=\frac{1}{2}LI^2\) & Magnetic-field storage. \\
Series resistors & \(R_{\mathrm{eq}}=\sum_k R_k\) & Same current through each resistor. \\
Parallel resistors & \(1/R_{\mathrm{eq}}=\sum_k 1/R_k\) & Same voltage across each branch. \\
Voltage divider & \(V_o=V_i R_2/(R_1+R_2)\) & Unloaded two-resistor divider. \\
Current divider & \(I_1=I\,R_2/(R_1+R_2)\) & Two parallel branches; current favors smaller resistance. \\
Thevenin/Norton & \(I_N=V_{Th}/R_{Th}\) & Linear one-port equivalents. \\
RMS sine & \(V_{\RMS}=V_p/\sqrt{2}\) & \(V_{pp}=2V_p\). \\
\bottomrule
\end{longtable}

\section{Impedance and Phasors}
Phasors represent sinusoidal steady-state signals. Unless stated otherwise,
phasor magnitudes are RMS.
\[
v(t)=\sqrt{2}\,\Re\{V\ee^{\jj\omega t}\},
\qquad
\frac{\dd}{\dd t}\longrightarrow \jj\omega .
\]
\[
Z_R=R,\qquad Z_L=\jj\omega L,\qquad
Z_C=\frac{1}{\jj\omega C}=-\frac{\jj}{\omega C}.
\]
\[
X_L=\omega L,\qquad X_C=\frac{1}{\omega C},\qquad
Z=R+\jj X,\qquad |Z|=\sqrt{R^2+X^2}.
\]
Use a quadrant-aware angle calculation:
\[
\theta=\operatorname{atan2}(X,R).
\]

\begin{mistakebox}
Do not mix reactance magnitude with signed impedance. For a capacitor,
\(X_C=1/(\omega C)>0\) is a magnitude, while \(Z_C=-\jj X_C\).
\end{mistakebox}

\section{Time Constants, Resonance, and Q}
\[
\tau_{RC}=RC,\qquad \tau_{RL}=\frac{L}{R}.
\]
\[
\omega_0=\frac{1}{\sqrt{LC}},
\qquad
f_0=\frac{1}{2\pi\sqrt{LC}}.
\]
For a \emph{series} RLC circuit,
\[
Z=R+\jj\left(\omega L-\frac{1}{\omega C}\right),
\qquad
Q_s=\frac{\omega_0 L}{R}=\frac{1}{\omega_0 R C}.
\]
Here \(Q_s\) is the series quality factor (sharpness of resonance). For many
narrowband resonant circuits the half-power bandwidth \(BW\) is
\[
BW\approx \frac{f_0}{Q}.
\]
Treat this as a useful approximation tied to half-power bandwidth, not as a
universal definition for every topology.

\section{Bode, Decibels, and Power}
\[
|G|_{\dB}=20\log_{10}|G(\jj\omega)|,
\qquad
\angle G=\arg G(\jj\omega).
\]
\[
10\log_{10}\frac{P_2}{P_1},\qquad
20\log_{10}\frac{V_2}{V_1}\quad \text{if impedance is unchanged.}
\]
One first-order pole contributes about \(-20~\dB/\mathrm{decade}\); one
first-order zero contributes about \(+20~\dB/\mathrm{decade}\).
\(\SI{20}{\decibel/decade}\) is about \(\SI{6}{\decibel/octave}\) (a decade is a
\(\times10\) frequency ratio; an octave a \(\times2\) ratio).

\begin{mistakebox}
\(-\SI{3}{\decibel}\) is half power and about \(0.707\) amplitude. It is not
half amplitude.
\end{mistakebox}

\section{s-Plane, Poles, and Feedback}
A pole \(s=\sigma+\jj\omega\) is a natural mode \(\ee^{\sigma t}\cos\omega t\):
left half-plane decays (stable), imaginary axis sustains, right half-plane grows.
The series RLC characteristic polynomial in standard form is
\[
s^2+2\zeta\omega_0 s+\omega_0^2,
\qquad
\alpha=\frac{R}{2L},
\qquad
\zeta=\frac{\alpha}{\omega_0}=\frac{1}{2Q},
\]
with underdamped poles \(s=-\alpha\pm\jj\omega_d\), \(\omega_d=\sqrt{\omega_0^2-\alpha^2}\).
Distance from the \(\jj\omega\) axis (\(\alpha\)) sets decay rate and bandwidth;
distance from the origin is \(\omega_0\). For a feedback loop (forward gain \(A\),
feedback fraction \(\beta\)),
\[
A_{\mathrm{cl}}=\frac{A}{1+A\beta},
\qquad
L=A\beta,
\qquad
\text{gain}\times\text{BW}=\text{const},
\]
and oscillation is \(L=-1\), i.e.\ \(|L|=1\), \(\angle L=-180^\circ\) (poles on the
\(\jj\omega\) axis).

\section{Filters, Matching, and Transformers}
\[
f_c=\frac{1}{2\pi RC}\quad \text{for a first-order RC cutoff.}
\]
An \(n\)-pole filter rolls off near \(-20n\,\si{\decibel}\) per decade. To match
real terminations \(R_{\text{hi}}>R_{\text{lo}}\) (the higher and lower resistances)
with an L-network,
\[
Q=\sqrt{\frac{R_{\text{hi}}}{R_{\text{lo}}}-1},
\qquad
X_{\text{series}}=Q\,R_{\text{lo}},
\qquad
X_{\text{shunt}}=\frac{R_{\text{hi}}}{Q}.
\]
\[
\frac{V_s}{V_p}=\frac{N_s}{N_p},
\qquad
\frac{I_s}{I_p}=\frac{N_p}{N_s},
\qquad
Z_p=\left(\frac{N_p}{N_s}\right)^2 Z_s.
\]
Matching networks cancel reactive impedance and transform the remaining
resistive part to the desired value. They can also restrict bandwidth, especially
when \(Q\) is high.

\section{Transmission Lines}
\[
\lambda=\frac{v_p}{f}=\frac{VF\,c}{f},
\qquad
\Gamma=\frac{Z_L-Z_0}{Z_L+Z_0}.
\]
\[
\mathrm{SWR}=\frac{1+|\Gamma|}{1-|\Gamma|},
\qquad
\mathrm{return\ loss}=-20\log_{10}|\Gamma|.
\]
Input impedance rotates along a lossless line,
\[
Z_{\mathrm{in}}(\ell)=Z_0\frac{Z_L+\jj Z_0\tan\beta\ell}{Z_0+\jj Z_L\tan\beta\ell},
\qquad
\beta=\frac{2\pi}{\lambda},
\]
giving \(Z_{\mathrm{in}}=Z_0^2/Z_L\) at a quarter wave and \(Z_{\mathrm{in}}=Z_L\)
at a half wave. Hence
\[
Z_t=\sqrt{Z_0Z_L}
\]
for a lossless quarter-wave transformer matching a real load \(Z_L\) to a line
with characteristic impedance \(Z_0\).

\section{What to Memorize and What to Derive}
Memorize the definitions of units, Ohm's law, \(Z_R\), \(Z_L\), \(Z_C\),
\(\omega=2\pi f\), \(f_0=1/(2\pi\sqrt{LC})\), decibel ratio rules, and the
transmission-line mismatch formulas. Derive dividers, time constants, and
series/parallel combinations from KCL, KVL, and equivalent impedance whenever
possible. That keeps the formula load small and reduces exam traps.

\begin{exambox}
For exam-style circuit questions, first ask: DC steady state or AC sinusoidal
steady state? Lumped circuit or transmission line? Series resonance or parallel
resonance? Power ratio or voltage ratio? The correct formula usually follows
from those choices.
\end{exambox}
